Long-horizon LLM agents are fundamentally bottlenecked by context: as interactions extend, useful information is overwhelmed by accumulated noise, undermining reasoning quality and operational efficiency. 
In this paper, we present GenericAgent (GA), a general-purpose agent system built around \textbf{context information density maximization}.
To operationalize this principle, GA integrates four components: (1) a minimal atomic tool set, (2) hierarchical on-demand memory, (3) reflection-driven self-evolution, and (4) structured browser extraction that filters raw web content into task-relevant signals.
These mechanisms continuously suppress low-value context and preserve high-signal information during long-horizon execution. 
Besides, GA further improves by turning execution traces into reusable procedures such as SOPs, code, and skills. 
We evaluate GA across task completion, token efficiency, tool use, memory, and web interaction. 
Experiments show strong performance and a consistently improved completion–cost trade-off compared with representative agent systems, with rapid convergence in cost and execution efficiency over repeated runs. 
We release GA as an open-source agent system at \href{https://github.com/lsdefine/GenericAgent/}{GitHub}.


